\section{Minimum Cost to Connect Sticks}
\label{sec:heap-connect-sticks}

\problemstatement{We are given $n$ sticks with given lengths. We may
repeatedly choose any two sticks, combining them into a single stick whose
length is the sum of the two, at a \textbf{cost equal to that sum}. We
repeat until only one stick remains. Find the \textbf{minimum total cost}
to combine all sticks into one.}

\begin{patternbox}
\emph{``Repeatedly combine the two smallest remaining items, at a cost
equal to their combination,''} is a direct heap-driven greedy: every
combination's cost depends on the lengths being combined, and combining
\textbf{small} sticks early keeps their (now-larger) sum from being
re-combined and re-charged many times over as the process continues. This
is structurally the same idea behind Huffman coding (building an optimal
prefix-free code by always merging the two least-frequent symbols first),
though no prior knowledge of Huffman coding is needed to derive it here.
\end{patternbox}

\subsection*{Understanding the problem carefully}

Every stick's original length is ultimately added into the cost once for
every combination step it participates in afterward, directly or as part
of a larger merged stick. A stick combined \emph{early} (when it is still
short) participates in fewer total dollars of cost across the remaining
steps than a stick combined \emph{late} (after it has already grown large
through earlier merges). So, to minimize total cost, we always want the
currently smallest available sticks to be the ones combined next --
deferring large sticks (or large merged sums) to participate in as few
further combination steps as possible.

\begin{ideabox}[Greedy Choice, Powered by a Min-Heap]
Push every stick length into a min-heap. Repeat until only one element
remains: pop the two smallest sticks, add their sum to the running total
cost, and push that sum back into the heap as a new (merged) stick.
\end{ideabox}

\subsection*{Why always combining the two smallest sticks is optimal}

This is an exchange argument identical in spirit to Huffman coding's
optimality proof: suppose an optimal combination sequence does not combine
the two globally smallest sticks $x \le y$ at some point, but instead
combines one of them with some larger stick $z$ first. One can show that
swapping the order -- combining $x$ and $y$ first instead -- never
increases the total cost, because $x$ and $y$ being the smallest means
they contribute the least possible amount per combination step they are
involved in, and delaying them only ever increases the number of
combination steps (each adding their value again, transitively, through
ever-larger merged sums) that their lengths end up contributing to.
Repeatedly applying this swap argument shows that always combining the two
currently smallest available sticks, at every step, is never worse than
any alternative ordering.

\subsection*{Algorithm}

\begin{enumerate}
  \item Push every stick length into a min-heap.
  \item Initialize $\textit{totalCost} \gets 0$.
  \item While the heap has more than one element:
  \begin{enumerate}
    \item Pop the two smallest values $a, b$.
    \item $\textit{totalCost} \gets \textit{totalCost} + (a+b)$.
    \item Push $a+b$ back onto the heap.
  \end{enumerate}
  \item Return \textit{totalCost}.
\end{enumerate}

\begin{algorithm}[H]
\caption{Minimum Cost to Connect Sticks}
\begin{algorithmic}[1]
\Procedure{ConnectSticks}{\textit{sticks}}
  \State $\textit{heap} \gets$ min-heap of all stick lengths
  \State $\textit{totalCost} \gets 0$
  \While{$|\textit{heap}| > 1$}
    \State $a \gets$ pop minimum, $b \gets$ pop minimum
    \State $\textit{totalCost} \gets \textit{totalCost} + a + b$
    \State push $a+b$ onto \textit{heap}
  \EndWhile
  \State \Return \textit{totalCost}
\EndProcedure
\end{algorithmic}
\end{algorithm}

\subsection*{Dry run}

\dryrun{Take $\textit{sticks} = [2,4,3]$.}

Heap $=\{2,3,4\}$.

\begin{itemize}
  \item Pop $2,3$. Cost $+=5$. Push $5$. Heap $=\{4,5\}$.
        $\textit{totalCost}=5$.
  \item Pop $4,5$. Cost $+=9$. Push $9$. Heap $=\{9\}$.
        $\textit{totalCost}=14$.
  \item Only one element remains; stop.
\end{itemize}

The minimum total cost is $14$.

\subsection*{C++ implementation}

\begin{lstlisting}
#include <bits/stdc++.h>
using namespace std;

int connectSticks(vector<int>& sticks) {
    priority_queue<int, vector<int>, greater<int>> minHeap(
        sticks.begin(), sticks.end());

    int totalCost = 0;
    while (minHeap.size() > 1) {
        int a = minHeap.top(); minHeap.pop();
        int b = minHeap.top(); minHeap.pop();

        totalCost += (a + b);
        minHeap.push(a + b);
    }
    return totalCost;
}

int main() {
    vector<int> sticks = {2, 4, 3};
    cout << "Minimum cost: " << connectSticks(sticks) << endl;
    return 0;
}
\end{lstlisting}

\begin{complexitybox}
\textbf{Time Complexity.} Building the initial heap costs $O(n)$ (using
the bottom-up build-heap technique from
Section~\ref{sec:heap-min-to-max}, which most standard library
implementations, including \texttt{std::priority\_queue}'s range
constructor, use internally). Each of the $n-1$ combination steps performs
two pops and one push, each $O(\log n)$, giving an overall time complexity
of
\[
  O(n \log n).
\]
\textbf{Space Complexity.} $O(n)$ for the heap.
\end{complexitybox}

\begin{takeawaybox}
``Repeatedly merge the two cheapest/smallest remaining items, paying a
cost proportional to the merge'' is the Huffman-coding greedy pattern,
recognizable on sight once you have seen it: a min-heap that always
surfaces the two cheapest current items, merged and pushed back, builds an
optimal merge order in $O(n \log n)$ without ever needing to consider
merge orders explicitly.
\end{takeawaybox}
